%% notes-tex/microbe-perturb-seq/sections/6-bottlenecks.tex
%% SOURCE: hand-authored synthesis. Every quantitative claim points back to the
%% section that derives it; no new numbers are introduced here.

\section{Bottlenecks and Development Priorities\secstatus{todo}}
\label{sec:bottlenecks}

What has to be true for 8--10 simultaneous perturbations to be measurable, and
what to do first given that we will start with single-guide CRISPRi and move to
two.

\subsection{Priority Ranking\secstatus{todo}}
\label{sec:ranked}

\pfstep{1. Guide capture (blocking)} Nothing else matters until guide identity is
recoverable from the RNA (Sec.~\ref{sec:guide-capture}). It has been treated here
as a construction problem rather than a cost problem, with a published yeast
template, and it may turn out not to be one: the 5$'$ assay captures guides with
a primer against the sgRNA scaffold rather than with an engineered capture
sequence, so an unmodified library might already be readable. Test that first,
because it costs one channel and it decides whether anything has to be built at
all. Let $q$ be the \emph{per-guide detection rate} -- the probability that
one guide present in a cell is actually seen in that cell's reads. A $k$-plex cell
is only usable if all $k$ of its guides are detected, so the usable fraction is
$q^{k}$, and the whole feasibility of multiplexing turns on which $q$ we inherit.
At Brandner et al.'s bacterial $q=0.25$ a 3-plex yields $0.25^{3}\approx1.6\%$ of
cells and an 8-plex is arithmetically hopeless. At the $q>0.71$ Nadal-Ribelles et
al.\ achieve in yeast, a 3-plex yields 36\% and even an 8-plex yields
6\%~\citep{nadal-ribellesSinglecellResolvedGenotypephenotype2025}. That is the
difference between multiplexing being impossible and merely expensive, and it is
one measurement. Raising $q$ is worth more than any other single change, because
it enters the multiplex ceiling exponentially while everything else enters
linearly -- and the pilot's first job is to find out which of those two regimes a
plasmid-borne guide barcode lands in.

\pfstep{2. mRNA fraction} At 5.75\% mRNA, 94 of every 100 reads buy ribosomal RNA
and the sequencing bill is inflated by roughly the reciprocal. Two independent
levers: shift the random-hexamer:oligo-dT ratio, since the 5$'$ coverage hexamers
buy is worthless for a 3$'$-counting screen; and add Cas9-guided rRNA depletion
after cDNA synthesis. Brandner et al.'s depletion took \org{E.\ coli} from 4.6\% to
60.2\% mRNA, but only 2.2\% to 10.4\% in \org{P.\ putida}, so the transfer to
yeast must be measured, not assumed. Worth $\sim$\$99\,k at genome scale
(Sec.~\ref{sec:budget}).

\pfstep{3. Cells per RT well} Brettner et al.\ verified 5{,}000 yeast cells per
well; microSPLiT runs 40{,}000--100{,}000 for bacteria. This number sets how many
cells one protocol run can barcode, and therefore how many runs a screen needs.
The 250-cells-per-gene split-pool row of \cref{tab:budgets} requires 13 protocol
runs at Brettner et al.'s loading; at 20{,}000 cells per well the same screen is
four runs, because a run would barcode $\sim$1.9\,M cells instead of
$\sim$480{,}000. That is $\sim$13 weeks of skilled bench work reduced to
$\sim$4, which is most of the labor argument for droplet
(Sec.~\ref{sec:budget}). It is also the cheapest experiment on this list. The
comparison it is measured against has moved: the droplet alternative is 18
preindexed channels rather than 137 plain ones, so the bench-time gap this
number closes is against 18, not 137.

\pfstep{4. Barcode space} A fourth ligation round is necessary at genome scale and
pays for itself within one run (Sec.~\ref{sec:collisions}). Its only cost is
molecular yield through an extra ligation, which the pilot must measure. Under
multiplexing the requirement tightens further, since a collision fabricates a gene
pair rather than merely diluting a main effect.

\pfstep{5. Transformation efficiency} The ceiling on multiplexing is library
construction, not measurement (Sec.~\ref{sec:library-ceiling}). Combinatorial
diversity generated by co-transforming several single-guide libraries, rather than
by constructing combinations, is the route worth costing -- it converts a
combinatorial cloning problem into a Poisson loading problem.

\pfstep{6. Effect size, now measured, and the two that took its place} This item
used to read that \cref{sec:effect-size} was unresolved and that everything in
\cref{sec:cells-per-pert} inherited its uncertainty. It is resolved: the median
responding gene moves 1.34-fold, which multiplied every depth-derived cell
requirement by $\sim$5.6 and is the largest single correction in the document.
What it did not resolve is the other two inputs to \cref{eq:design}, and both are
now the thing to measure. $\varphi_j$, gene-level overdispersion, sets the floor
by itself and is estimable from any yeast count matrix. $p_j$, the target gene's
transcriptome share, is currently assembled from two sources that do not belong
together and is worth a five-fold spread in cells per perturbation
(\cref{sec:transcript-content}). Both come out of the first pilot's count matrix
at no extra cost, which is what makes them cheaper than anything above them on
this list and is why they are ranked here rather than as an experiment.

\pfstep{7. Preindexed droplet, as a hedge rather than a plan} The item that would
most change the ranking above is the one nothing on this list depends on.
Preindexing takes the droplet route from $6.3\times$ the cost of depleted
split-pool to $1.9\times$ (Sec.~\ref{sec:budget}), which converts droplet from a
platform we have ruled out on price into a live fallback if items 1 or 3 fail.
It ranks last because it is the least evidenced: items 1--6 are quantities to be
measured on a chemistry that has been run in yeast, whereas this is a chemistry
that has never been run in any microbe, and the same in-cell reverse
transcription it needs is what item 3 is already testing. That ordering is worth
stating rather than leaving implicit, because the cost figure is attractive
enough to invite reordering it, and the figure and the evidence run opposite
ways.

\subsection{The Host Organism, and Why the First Experiment Carries No
Perturbations\secstatus{todo}}
\label{sec:host}

Everything above is written for \org{S.\ cerevisiae}. The intended production
host is \org{Pichia kudriavzevii}, the same species as \org{Issatchenkia
orientalis} and as the clinical isolate \org{Candida krusei}, whose genomes are
``collinear \dots\ 99.6\% identical in DNA
sequence''~\citep{douglassPopulationGenomicsCandida2018}. Moving there changes
the risk profile sharply, and in a direction that decides what to run first.

\notesec{What exists} A complete reference genome, five chromosomes and
10.8\,Mb with 5{,}139 protein-coding genes annotated, so read alignment and
guide design are not obstacles. Catalytic Cas9 works well: an \gene{ADE2}
disruption efficiency of 97.0\% using an \gene{RPR1}$'$-tRNA$^{\text{Leu}}$
promoter for the sgRNA, and multiplexed deletions at 72.8\% to 89.9\% for pairs
and 46.7\% for a triple~\citep{tranDevelopmentCRISPRCas9Based2019}. A landing-pad
system reaches greater than 80\% integration efficiency at 27 characterized loci
and supports simultaneous integration of two to five
genes~\citep{fatmaLandingPadSystem2023}.

\notesec{What is not published, which is not the same as what does not exist}
Two of the gaps below are closed in house and the distinction matters, because a
capability that exists but is unpublished is a validation task rather than a
research project. There is no published dCas9 CRISPRi or CRISPRa system in this
organism -- every CRISPR result above is catalytic knockout or integration --
but a working CRISPRi system and an assembled genome from JGI are both held
here\external{in-house, unpublished}. Read against that, the missing piece is not
the perturbation tool; it is everything downstream of it.

What remains genuinely open is delivery and scale. Episomal plasmids are
mitotically unstable, with about 55\% of cells still expressing a plasmid-borne
reporter after five days~\citep{tranDevelopmentCRISPRCas9Based2019}, which is a
problem for the delivery route of \cref{sec:delivery} specifically; a
centromere-like sequence found in the same organism improves plasmid stability
and is the obvious thing to build
on~\citep{caoGeneticToolboxMetabolic2020}. And the only
transformation-efficiency anchor available, roughly $10^{3}$ colony-forming
units per microgram, is three to five orders of magnitude below what
\org{S.\ cerevisiae} achieves. \Cref{sec:library-ceiling} showed that library
scale is set by transformation; if that figure is representative, pooled
library-scale screening in this host is not currently reachable and the gap is
not one a budget closes. The published strain used by the bioproduction
community is also diploid, which changes what a single guide does and what a
knockout means.

\begin{keybox}
\textbf{The first experiment carries no perturbations, and the reason is footing
rather than a missing tool.} Nobody on this project has run single-cell RNA-seq
on a walled microbe, so the first run is worth spending on learning the assay:
whether the cell wall yields to the in-droplet zymolyase of \cref{sec:jariani} in
this species, how many mRNA UMIs per cell are recovered, and what fraction of
reads is ribosomal. That is the collaborator's recommendation and it is also what
the arithmetic supports, since the third of those is the largest cost lever in
the document (\cref{sec:rrna}) and is a proportion, answerable at any sequencing
depth. None of the three needs a library, so none of them waits on the delivery
and scale questions above. If a run with no perturbations looks like a thin
return for a channel, the way to thicken it is to spread the same run across
several environments (\cref{sec:environments}), which costs barcode capacity that
is otherwise idle rather than costing genetics that do not exist. Running this
first also frees the guide-capture question of \cref{sec:guide-capture} to be
settled in \org{S.\ cerevisiae}, where the library already exists, in parallel
rather than in sequence.
\end{keybox}

\notesec{Two runs, not one: a MiSeq to check, a NovaSeq lane to profile} The
core's own advice is to use the MiSeq for quality control rather than for
production, and that maps cleanly onto this experiment because the two things
being asked are different questions.

The check is whether the chemistry worked at all: did the wall yield, did cells
barcode, what fraction of reads is ribosomal. All three are proportions, so they
are answerable at almost any depth. A MiSeq i100 25\,M paired-end run
(\$1{,}450) carries $\sim$1{,}250 cells at 10x's recommended 20{,}000 read pairs
per cell, and the \$398 titration run answers the barcode question alone. Either
is cheap enough to repeat if the first attempt fails, which is the property that
matters at this stage: \cref{sec:staging} notes that cDNA traces and library
concentration do not predict barcoding success, so there is no earlier
checkpoint than sequencing something.

The profiling run is a different job. A first transcriptome atlas of this
organism is a dataset we would want to keep and analyze, and 1{,}250 cells is
not one. \Cref{sec:sequencing} showed that the core's gap sits exactly here, but
the split-lane row is the way through it: a NovaSeq X 10B split lane gives
600\,M read pairs for \$1{,}160, which is $\sim$30{,}000 cells at full depth,
for less than the MiSeq run costs. So the sequencing is not what constrains this
experiment; the 10x channel is, at 20{,}000 cells recovered.

\begin{keybox}
\textbf{Sequence the pilot twice, and the second time is cheaper than the first.}
Run the MiSeq to confirm the chemistry, then a NovaSeq X 10B split lane
(\$1{,}160 for 600\,M read pairs) for the actual profile. The split lane costs
less than the 25\,M MiSeq run and delivers 24 times the reads, so the only
reason to sequence shallowly here is turnaround, not money. Depth is not the
binding constraint on this experiment and neither is cost.
\end{keybox}

\subsection{Staged Execution Plan\secstatus{todo}}
\label{sec:staging}

%% TODO: firm up once the Kemmeren/Sameith effect-size analysis lands and the
%% guide-capture construct is chosen. Sketch only.

\begin{enumerate}[leftmargin=1.4em,itemsep=2pt,topsep=2pt]
  \item \textbf{Unperturbed \org{P.\ kudriavzevii}} (\cref{sec:host}). No
    library and no guides, so it depends on none of the delivery and scale
    questions still open in this host, and it is where the assay itself gets
    learned. It returns the wall-digestion verdict, the mRNA UMIs per cell, and
    the ribosomal fraction, which is the largest cost lever in the document.
    Sequence it twice: a MiSeq run to confirm the chemistry, then a NovaSeq X 10B
    split lane for a profile worth keeping. If the channel looks under-used
    carrying one condition, run it across several environments instead of adding
    guides. This is the current first experiment.
  \item \textbf{Test 5$'$ guide capture on the existing
    \org{S.\ cerevisiae} library}, in parallel with step 1 rather than after it,
    because it uses a different organism and a library that already exists. One
    channel on the core's GEM-X 5$'$ kit answers whether the scaffold primer
    reverse transcribes an unmodified Pol III guide (\cref{sec:guide-capture}).
    If it does, the construction project below is not needed.
  \item \textbf{Construct and validate the poly(A) guide barcode}, \emph{only if
    step 2 fails}, in an arrayed set of $\sim$10 strains with known knockdown
    phenotypes, scored against bulk RNA-seq from the same strains. Ground truth
    is only available while the design is arrayed; a pool cannot tell you what
    was in a given cell.
  \item \textbf{Establish the chemistry numbers} the model is most sensitive to:
    mRNA fraction after depletion, cells per RT well, guide-assignment rate $q$,
    the read-depth saturation curve, and -- from the same data, at no extra cost -- the transcriptome share $p_j$ of a few representative target genes, which
    \cref{sec:transcript-content} shows is currently worth a five-fold spread in
    cells per perturbation and is the largest unforced uncertainty in the budget. One \$398 MiSeq titration run protects a
    \$3{,}180 NovaSeq lane, and there is no earlier checkpoint -- Brettner et al. is
    explicit that cDNA traces and library concentration do not predict barcoding
    success.
  \item \textbf{Answer the preindexing question from data step 4 already
    produces}, at no additional cost. Whether a fixed, wall-digested yeast cell
    reverse transcribes in situ and then survives handling is measured by the
    round-1 step of the split-pool pilot; what step 4 does not answer is whether
    such a cell survives the pressures of a microfluidic chip, which is one
    channel to test and is the only part of Sec.~\ref{sec:scifi} with no yeast
    precedent at all. Worth running only if step 1 or step 3 goes badly, because
    it buys a fallback rather than advancing the main route.
  \item \textbf{Single-guide pilot} on a focused sub-library, at the 100-cell
    floor.
  \item \textbf{Two-guide screen} on the focused metabolic sub-library, where
    named pairs are reachable, to measure how guide detection and collisions
    actually degrade at $k=2$ before trusting any extrapolation to $k=8$.
  \item \textbf{Genome-scale single-guide screen}, four barcoding rounds, one
    environment.
\end{enumerate}

\subsection{Outlook\secstatus{todo}}
\label{sec:vision}

%% TODO: this subsection is a placeholder for the longer-horizon argument and is
%% deliberately thin -- it is the part that becomes a thesis chapter rather than a
%% method plan.

Two extensions change what the platform is for, and both are worth designing
toward now rather than retrofitting.

\notesec{Splitting the sample -- and why the obvious version does not work}
Pairing a transcriptome with mass spectrometry on the \emph{same} perturbation is
the extension that would matter most for metabolic engineering, since flux is what
we actually want and it cannot be inferred from expression alone.
Fulcher et al.\ \citep{fulcherParallelMeasurementTranscriptomes2024} is the precedent for parallel
transcriptome and proteome from the same single cells. But the naive version -- take
half the contents of a compartment -- is not available in either format, and for
different reasons worth separating.

A semipermeable capsule cannot be aliquoted. It is a hydrogel shell with a pore
cutoff around 30\,nm: enzymes up to $\sim$160\,kDa diffuse in, and that
permeability is the entire point, since it is what lets reagents reach the cells
without breaking the compartment~\citep{baronasHighthroughputSinglecellOmics2026}.
There is no valve; the pores are a property of the material, not a gate that can
be closed. And the same cutoff is fatal to the specific application: \emph{small
molecules leak}. Metabolites are precisely the species that pass a 30\,nm pore
freely, so an SPC cannot hold a metabolome in the first place, let alone divide
one. Proteins are the more plausible target, since much of a proteome sits above
the cutoff -- but a metabolomic readout from SPCs is not a matter of engineering
the split, it is excluded by the format.

Two routes remain, and they split the population rather than the compartment.
Capsules are FACS-sortable and their cells can be released alive by brief protease
treatment, so whole capsules can be \emph{routed} -- one fraction sequenced,
another pooled for bulk mass spectrometry -- which gives paired measurements over
matched sub-populations rather than over identical cells. Alternatively, because
each capsule holds a \emph{clonal} population expanded from one founder, releasing
that clone and dividing the cells does give two aliquots of the same genotype;
what is lost is the capsule identity, which then has to be re-established by a
barcode rather than by physical containment.

\notesec{The unsolved part} Either route needs the mass-spectrometry channel to
carry perturbation identity, and MS has no equivalent of a sequencing barcode. The
plausible designs are chemical isobaric labeling of routed fractions, or
physically arraying sorted capsules into wells so that position encodes identity -- which trades throughput for it, and is the obvious tension to think about
before this becomes a plan. Isotope labeling for flux sits downstream of solving
that.

\notesec{Environments as a second axis} Treated in full in
\cref{sec:environments}, which is where the arithmetic sits: environments
multiply cell demand linearly where named gene combinations grow quadratically,
and split-pool's round-1 plate already has the barcode capacity to carry them.
